% this is a test for latex Packages
% Author: Fauzy Madani <fauzy_madani@smknegeri1garut.sch.id>
% This style and looks is inspired from a reddit post:
% https://www.reddit.com/r/LaTeX/comments/1klt7wd/help_latex_chapter_style/#lightbox
\documentclass{book}

% ========== File contents untuk Zallman.sty ==========
\usepackage{filecontents}
\begin{filecontents*}{Zallman.sty}
\NeedsTeXFormat{LaTeX2e}
\ProvidesPackage{Zallman}[2007/11/24 v1.0 Zallman CFR]

\input Zallman.fd

\DeclareRobustCommand{\Zallmanfamily}{
  \fontencoding{U}%
  \fontseries{xl}%
  \fontshape{n}%
  \fontfamily{Zallman}%
  \selectfont}
\DeclareTextFontCommand{\zall}{\Zallmanfamily}
\endinput
\end{filecontents*}

% ========== Packages ==========
\usepackage[a4paper, top=2cm, bottom=3cm, left=2.5cm, right=2.5cm]{geometry}
\usepackage{titlesec}
\usepackage{tikz}
\usepackage{array}
\usepackage{xcolor}
\usepackage{etoolbox}
\usepackage[hidelinks]{hyperref} % hyperlink for TOC
\usepackage{lipsum} % dummy text
\usepackage{lettrine} % for drop cap

% ========== Load Zallman ==========
\usepackage{Zallman}
\definecolor{Maroon}{rgb}{0.5, 0, 0}
\renewcommand{\LettrineFontHook}{\color{Maroon}\Zallmanfamily}

% ========== Chapter spacing tweak ==========
\makeatletter
\patchcmd{\@makechapterhead}
  {\vspace*{50\p@}}
  {\vspace*{10pt}}
  {}{}
\makeatother

% ========== Colors ==========
\definecolor{deepred}{RGB}{153, 0, 0}

% ========== Chapter style ==========
\titleformat{\chapter}[display]
  {\normalfont\Large\bfseries}
  {%
    \begin{tikzpicture}
      \fill[deepred] (0,0) rectangle (2,-2.8);
      \node[text=white, font=\small\bfseries, anchor=base] at (1,-0.3) {CHAPTER};
      \node[text=white, font=\fontsize{30}{36}\selectfont\bfseries] at (1,-1.5) {\thechapter};
    \end{tikzpicture}
  }
  {0.5em}
  {\MakeUppercase}
  [
    \vspace{0.3em}%
    \small\itshape\mdseries\raggedleft
    "The true logic of this world is in the calculus of probabilities." \\
    \rule{\linewidth}{0.4pt} \\
    \textbf{--- James C. Maxwell}
  ]

\titlespacing*{\chapter}{0pt}{0pt}{2em}

% ========== Section & Subsection style ==========
\titleformat{\section}
  {\normalfont\bfseries}
  {\thesection}
  {1em}
  {}

\titleformat{\subsection}
  {\normalfont\bfseries}
  {\thesubsection}
  {1em}
  {}

% ========== Column types ==========
\newcolumntype{L}[1]{>{\raggedright\arraybackslash}p{#1}}
\newcolumntype{R}[1]{>{\raggedleft\arraybackslash}p{#1}}

% ========== Document starts ==========
\begin{document}

% === Table of Contents ===
\tableofcontents
\clearpage

% === CHAPTER 1 ===
\chapter{How to Use This Book}

\section{Introduction}
\lettrine{F}{ irst of all}, welcome to the world of Deep Learning Interviews. \lipsum[1]

\subsection{What makes this book so valuable}
\lettrine{D}{ eciphering} dead languages. Detecting malignant tumours. Predicting natural disasters. \lipsum[2]

% === CHAPTER 2 ===
\chapter{Deep Learning Basics}

\section{What is a Neural Network?}
\lettrine{A}{ rtificial} neurons are the building blocks. \lipsum[3]

\subsection{Perceptron}
\lettrine{T}{ he} perceptron is a binary classifier. \lipsum[4]

\section{Training Deep Networks}
\subsection{Backpropagation}
\lettrine{B}{ ackpropagation} is an essential technique. \lipsum[5]

% === CHAPTER 3 ===
\chapter{Advanced Topics}

\section{GANs and VAEs}
\lettrine{G}{ enerative} models like GANs are powerful. \lipsum[6]

% === CHAPTER 4 ===
\chapter{Interview Questions}

\section{Theory}
\lettrine{T}{ his} section covers theoretical concepts. \lipsum[7]

\section{Coding}
\lettrine{I}{ nterviewers} often ask you to write code. \lipsum[8]

\end{document}

